\section{Layer~1 Results}\label{sec:fraudres}
Among classical models, random forest reached PR-AUC \ClsRfPr{} and XGBoost \ClsXgbPr{} (logistic regression \ClsLrPr{}; isolation forest \ClsIfPr{}). XGBoost was retained for production because of much lower latency at a similar PR-AUC, not because of the highest point estimate.

In the decisive quantum experiment (\LsSeeds{} paired seeds), the local-sum observable improved PR-AUC over the global observable by a mean of \LsDelta{} (bootstrap 95\% CI [\LsCIlo{}, \LsCIhi{}]; one-sided Wilcoxon $p=\LsWilcoxon{}$; permutation $p=\LsPerm{}$; local-sum better in \LsWins{} of \LsSeeds{} seeds). Fig.~\ref{fig:paired} shows the paired per-seed comparison. This supports an architecture-level effect of observable choice. It does \emph{not} establish a system-level quantum advantage: on the locked confirmation split, XGBoost alone scored \CapXgb{}, XGBoost plus the local-sum QNN \CapXgbLocal{}, and XGBoost plus a matched classical control \CapXgbControl{}; the bootstrap CI of the quantum increment over XGBoost alone, [\CapCIlo{}, \CapCIhi{}], includes zero. Every other quantum configuration evaluated (QSVM, VQC, Chow--Liu circuits, QBoost, hybrid fusion architectures, quantum-kernel anomaly detection at scale) likewise failed to beat its classical counterpart; these negative results are retained in the repository.

\begin{figure}[t]
\centering
\includegraphics[width=\figw]{paired_seed_deltas}
\caption{Paired PR-AUC comparison of the global and local-sum QNN across matched seeds.}
\label{fig:paired}
\end{figure}
